% 建模活动单 · 第六组
% 两个活动都用基本不等式处理“约束下的最值”：先锁死取等条件，再代数求解。
\section{活动十二：方盒怎样最省料——基本不等式}

\subsection{情境}
纸盒厂要生产一种容积为 $4000\ \mathrm{cm^3}$ 的长方体包装盒。
纸板成本近似与盒子的表面积成正比，老板希望“同样装 $4000\ \mathrm{cm^3}$，用的纸板越少越好”。

\subsection{提出问题}
长、宽、高各取多少时，盒子的表面积最小？最小是多少？

\subsection{选模型}
设长、宽、高为 $x,y,z$，则约束与目标分别为
\[
  xyz=V\quad(V=4000),\qquad S=2(xy+yz+zx).
\]
“乘积固定，求和的最小值”正是基本不等式的用武之地。

\subsection{求解}
由三元基本不等式
\[
  xy+yz+zx\ge 3\sqrt[3]{(xy)(yz)(zx)}=3\sqrt[3]{(xyz)^2}=3\sqrt[3]{V^2},
\]
当且仅当 $xy=yz=zx$，即 $x=y=z$ 时取等号。于是
\[
  S\ge 6\sqrt[3]{V^2}=6\sqrt[3]{4000^2}\approx 1511.9\ (\mathrm{cm^2}),
\]
最优尺寸为立方体
\[
  a=\sqrt[3]{4000}=10\sqrt[3]{4}\approx 15.87\ (\mathrm{cm}).
\]

\subsection{与常见方案的比较}
\begin{center}
\footnotesize
\setlength{\tabcolsep}{6pt}
\renewcommand{\arraystretch}{1.4}
\begin{tabular}{p{3.2cm}p{4.6cm}p{3.0cm}p{2.4cm}}
\toprule
方案 & 尺寸（cm） & 表面积（$\mathrm{cm^2}$） & 与最优之差 \\
\midrule
立方体（最优） & $15.87\times15.87\times15.87$ & $1511.9$ & — \\
$1:2:1$ & $12.60\times25.20\times12.60$ & $1587.4$ & $+5.0\%$ \\
$2:2:1$ & $20\times20\times10$ & $1600$ & $+5.8\%$ \\
\bottomrule
\end{tabular}
\end{center}
$1:2:1$ 方案中，由 $x\cdot 2x\cdot x=4000$ 得 $x=\sqrt[3]{2000}\approx12.60$，
$S=10x^2\approx1587.4$；$20\times20\times10$ 的体积虽然正好 $4000$，
但形状偏离立方体更远。可见“体积对”不等于“省料”，偏离最优形状 $1:2:1$
就多花约 $5\%$ 的纸板。

\subsection{检验与反思}
用拉格朗日乘数法看（选学）：$\nabla S=\lambda\nabla(xyz)$ 给出
$y+z=\lambda yz,\ z+x=\lambda zx,\ x+y=\lambda xy$，化简后同样是 $x=y=z$。
两种方法指向同一结论——最优点处三个“面面积”$xy,yz,zx$ 相等。
\key{“乘积固定求和最小”，先用基本不等式锁定取等条件（相等），再代回约束求尺寸。}

\section{活动十三：易拉罐的高矮胖瘦——基本不等式}

\subsection{情境}
某厂要设计一款 $355\ \mathrm{mL}$ 的圆柱形易拉罐。铝材用量正比于表面积，
从省料角度看，罐身高与直径的比值 $h:d$ 取多少最好？

\subsection{提出问题}
怎样的 $r,h$ 使表面积最小？这时 $h$ 与 $d$ 有什么关系？

\subsection{选模型与求解}
容积固定：$\pi r^2h=V=355$。表面积
\[
  S=2\pi r^2+2\pi rh=2\pi r^2+\frac{2V}{r}
\]
把 $\dfrac{2V}{r}$ 拆成两项，用三元基本不等式：
\[
  S=2\pi r^2+\frac{V}{r}+\frac{V}{r}
   \ge 3\sqrt[3]{2\pi r^2\cdot\frac{V}{r}\cdot\frac{V}{r}}
   =3\sqrt[3]{2\pi V^2}\approx 277.5\ (\mathrm{cm^2}),
\]
取等条件为 $2\pi r^2=\dfrac{V}{r}$，即 $r^3=\dfrac{V}{2\pi}$。代回 $h=\dfrac{V}{\pi r^2}$
得 $h=2r=d$：\emph{高等于直径}时最省料，此时
\[
  r\approx 3.84\ \mathrm{cm},\quad d=h\approx 7.67\ \mathrm{cm},\quad S_{\min}\approx 277.5\ \mathrm{cm^2}.
\]

\subsection{若比例自定：$h=2d$}
若按外观需要先定 $h=2d$，则 $V=\pi\left(\dfrac d2\right)^2\cdot 2d=\dfrac{\pi d^3}{2}$，
\[
  d=\sqrt[3]{\frac{2V}{\pi}}\approx 6.09\ \mathrm{cm},\quad h\approx 12.18\ \mathrm{cm},
  \quad S=\frac52\pi d^2\approx 291.4\ \mathrm{cm^2},
\]
比数学最优多约 $5.0\%$——和自由优化相比，“拍脑袋定比例”并不便宜多少。

\subsection{检验与反思：为什么市售易拉罐不按数学最优}
市售 $355\ \mathrm{mL}$ 易拉罐多为 $d\approx6.6\ \mathrm{cm}$、$h\approx12.2\ \mathrm{cm}$
（$h:d\approx1.85$，罐内预留顶隙，实际容积约 $417\ \mathrm{mL}$）。
若保持这一比例、按 $355\ \mathrm{mL}$ 等比折算，表面积约 $288.5\ \mathrm{cm^2}$，
仍比数学最优多约 $4\%$；顶底局部加厚还会进一步拉大差距。原因不在数学，而在工程与使用：
\begin{itemize}
  \item 顶盖强度：罐内承压，顶面越大越易变形，矮粗罐必须把顶盖加厚，反而更费料；
  \item 堆叠与运输：$h:d\approx2$ 的重心与码放层数更稳，托盘装卸不易倒罐；
  \item 手感：单手握持、入口角度都偏好细高外形；
  \item 货架效果：竖立陈列时细高罐的正面展示面积更大。
\end{itemize}
修正模型思路：把人的需求写成\emph{约束或惩罚项}，而不是改目标函数。例如规定
$h\ge 1.5\,d$（强度与握持下限），则最优解被约束推到边界 $h=1.5\,d$ 处取得，
由 $V=\dfrac{3\pi d^3}{8}$ 得
\[
  d\approx 6.70\ \mathrm{cm},\quad h\approx 10.06\ \mathrm{cm},\quad S\approx 282.4\ \mathrm{cm^2},
\]
只比无约束最优多 $1.8\%$。也可以改成加权目标 $J=S+\lambda\left(\dfrac hd-\rho_0\right)^2$，
用 $\lambda$ 权衡材料成本与体验、强度。
\key{工程最优常不在数学极值点，而在约束边界上；建模时先把“人”的需求翻译成约束或惩罚项。}
